\section{Discussion}
\label{sec:discussion}
\subsection{Interpretation}
The protocol checks and onboarding replay address different steps in device
integration. Protocol connectivity demonstrates communication with the tested
simulator endpoints, while the replay checks whether one configuration exposes
its declared acquisition and control nodes. Neither establishes that an
arbitrary device can be integrated without adaptation. The entry-path checks
show which records are available for each tested call; separate environments
and absent handler-level traces prevent a claim of route equivalence or direct
verification of a shared handler. In the separate approval-context test, stale product and run identities
reached mock writes when the tuple comparison was removed but were rejected
when it was enabled. This contrast isolates the check's contribution within
that test setup.

The four missions make the distinction between a control request and a
process outcome concrete. In the molding and film examples, several requests
precede the final in-band observation; in annealing, the record also contains
a reverted candidate. In wastewater treatment, the aeration trajectory and
the later mission summary refer to different observations. Keeping these
records distinct permits the sequence to be inspected without treating a
command as proof of a desired process response.

The missions also illustrate the role of scene-specific configuration.
Part weight, film thickness, hardness, and dissolved oxygen have different
units, target bands, and available actuators. The operation--response
structure accommodates these differences, but the trajectories do not
identify which architectural component caused an outcome. Related approaches
likewise connect industrial operations to semantic skills, equipment context,
or process validation \cite{koehle2025semanticrecipes,graul2026aas,schall2026safe}.
The present study complements those designs by describing a shared operation
path and showing recorded responses in four simulated scenes; the studies use
different tasks and do not provide a matched performance comparison.

\subsection{Scope and Limitations}
The optimization evidence comprises one deterministic simulated trajectory per scene,
without independent replicates or a matched alternative controller. It cannot
establish statistical repeatability, superiority to another platform, or the
effect of individual controls. The missions use direct-node operations and
do not test the full range of entry routes or language-model decisions
described by the architecture. Final tensile-strength and surface-defect
readings for annealing and the complete final wastewater effluent vector are
missing, limiting conclusions about secondary constraints. The process model
and simulator configuration also limit transfer to physical production.

The mission records show automated judgments, not human pre-approval. The
approval-context test uses simulated decisions and mock drivers; it tests a
software condition, not operator performance or the effectiveness of human
review \cite{parasuramanmanzey2010complacency,bainbridge1983ironies}.
Software operation restrictions are not a complete security or safety
assessment \cite{greshake2023injection,liu2024pinjection,iec62443-3-3};
independent plant protection remains necessary \cite{iec61511}. Future work
should repeat missions across initial conditions, preserve all terminal
quality measurements, compare matched controllers and entry paths, and
evaluate the system with suitable physical equipment and operator oversight.

These results separate three integration decisions: which entry a caller uses,
which scene mapping an operation resolves, and which evidence an owning service
retains. This separation lets an application change an agent engine or a device
mapping without redefining the operation record. The evidence supports this as
a software structure; it does not measure engineering effort, operating cost,
or reliability gains, which would require matched deployments.

\section{Conclusion}
\label{sec:conclusion}
AgentWorkShop connects industrial task context and scene-specific device
operations through shared dispatch, configured drivers, and separately
retained command and observation records. Simulator connectivity and one
onboarding replay exercise configured protocol and node access; separate
entry-path checks record route-specific evidence, and a mock test checks stale
approval-context rejection. Four local simulation missions
illustrate the operation--response sequence across different process
objectives; their reported final primary values lie within their respective
target bands. The examples show how process evidence can be interpreted alongside requested
actions. The single deterministic trajectories and missing secondary
measurements leave repeatability, comparative benefit, human-review
effectiveness, and physical-plant performance for future evaluation.

\noindent\textbf{Data Availability:} The historical Supplementary File~S1
contains route and integration-check records and the onboarding replay log;
it does not contain the current four-mission traces. A separate supplementary
evidence bundle contains pseudonymized mission logs and selected
registered-harness records. The manuscript source, vector figures, and
figure-source data accompany the submission package. No physical-plant data
are included.
